\par\Needspace{8\baselineskip}\bigskip
\pdfbookmark[3]{Darstellung, 2}{pallas-II-article2}
\begin{center}2.\end{center}
\noindent\hspace*{\parindent}Setzt man in Gleichung 16 \(\mathrm{d}x^2+\mathrm{d}y^2+\mathrm{d}z^2=\mathrm{d}r^2+\frac{pm\mathrm{d}t^2}{rr}\mathpunct{\raisebox{.45ex}{,}}\) so erhält man \(\frac{\mathrm{d}r^2}{\mathrm{d}t^2}=-\frac{m}{a}+\frac{2m}{r}-\frac{pm}{rr}\mathpunct{\raisebox{.45ex}{;}}\) an den Punkten der Bahn, an denen \(r\) ein Maximum \SourceLeaf{449}oder ein Minimum annimmt, muss daher \(\frac{m}{a}-\frac{2m}{r}+\frac{pm}{rr}=0\) gelten. Das ergibt für \(r\) zwei Werte, nämlich \(r=a+\BareRoot{(aa-ap)}\) und \(r=a-\BareRoot{(aa-ap)}\). Daraus ist ersichtlich, dass \(2a\) die große Achse ist und \(\frac{\BareRoot{(aa-ap)}}{a}\mathpunct{\raisebox{.45ex}{,}}\) wobei wir diese Größe \(=e\) setzen, die Exzentrizität ist. Also gilt
\[
\begin{gathered}
p=a(1-ee)\\[4pt]
\mathrm{d}t^2=\frac{arr\mathrm{d}r^2}{m(aaee-(a-r)^2)}\mathpunct{\cdot}
\end{gathered}
\]
\par Um diese Gleichung zu integrieren, setzen wir
\[
\begin{array}{rl}17)\quad&a-r=ae\cos E,\end{array}
\]
woraus folgt
\[
\mathrm{d}t^2=\frac{arr\mathrm{d}E^2}{m}\mathpunct{\raisebox{.45ex}{;}}
\]
Setzt man nun, was zulässig ist, voraus, dass \(E\) stets zunimmt, so gilt
\[
\mathrm{d}t=\BareRoot{\frac{a}{m}}\cdot r\mathrm{d}E=\BareRoot{\frac{a^3}{m}}\cdot(1-e\cos E)\mathrm{d}E.
\]
\par Setzen wir also
\[
\begin{array}{rl}
18)\quad&\BareRoot{\dfrac{m}{a^3}}=n\\[6pt]
19)\quad&E-e\sin E-nt=\epsilon,
\end{array}
\]
so zeigt die Gleichung \(n\mathrm{d}t=(1-e\cos E)\mathrm{d}E\), dass \(\epsilon\) eine Konstante ist. Man sieht, dass \(E\) das bezeichnet, was die Astronomen die exzentrische Anomalie nennen, \(nt+\epsilon\) die mittlere Anomalie, \(n\) die mittlere Bewegung in der Zeiteinheit und \(\epsilon\) die Epoche der mittleren Anomalie zur Zeit null. Außerdem ist klar, dass \(E\) in der Zeit eines Umlaufs um \(360^{0}\) oder um \(2\pi\) zunehmen muss, wobei \(2\pi\) der Umfang eines Kreises vom Radius \(1\) ist. Bezeichnet man diese Zeit also mit \(T\), so gilt
\[
nT=2\pi\qquad\text{oder}\qquad T=2\pi\BareRoot{\frac{a^3}{m}}\mathpunct{\raisebox{.45ex}{,}}
\]
was auf den Satz von \textsc{Kepler} hinausläuft.
\par Damit erhält man auch das Maß der Massen, sobald man die Einheiten der Zeit und des Raumes gewählt hat, nämlich
\[
m=\frac{4\pi\pi a^3}{TT}\mathpunct{\cdot}
\]
\par Wählt man also als die eine Einheit den mittleren Sonnentag und als die andere den \SourceLeaf{450}Wert von \(a\) in der Erdbahn und setzt außerdem \(T=365{,}2[563835]\) voraus, so erhält man
\begin{center}
\textit{Summe der Massen der Sonne und der Erde}\par\nobreak
\mbox{. . . . \(0{,}000295913[0424]\)}
\end{center}
oder, wenn man annimmt, dass die erstere \(354710\)-mal so groß ist wie die letztere
\begin{center}
\textit{Masse der Sonne} . . . . \(0{,}000295912[2082]\),
\end{center}
deren gewöhnlicher Logarithmus \(6{,}471\;1628\;828-10\) ist.
